\chapter{Urinary Incontinence}
\index{Urinary Incontinence}

\section*{Definition and Overview}
Urinary incontinence is the involuntary leakage of urine. It is a common and distressing problem that affects people of all ages, though it becomes more frequent with age and after childbirth. The two most common types are stress incontinence, in which urine leaks with coughing, sneezing, laughing, or exercise because of weakness of the pelvic floor muscles and sphincter, and urge incontinence (overactive bladder), in which a sudden strong urge to pass urine is followed by leakage before reaching the toilet. Mixed incontinence combines both. Many people wrongly believe that leakage is a normal part of ageing or childbirth and that nothing can be done. In fact, incontinence can almost always be improved and often cured with the right assessment and treatment.

\section*{Epidemiology}
Urinary incontinence affects around one in three women and up to one in ten men, and it becomes more common with age. Stress incontinence is the most common type in younger and middle-aged women, while urge and mixed incontinence increase with age in both sexes. Around a third of women leak urine in pregnancy, and vaginal birth, especially with a large baby or instrumental delivery, increases the risk. Men most often develop incontinence after prostate surgery or with prostate disease. Despite its frequency, only a minority of affected people seek help, often because of embarrassment or the mistaken belief that nothing can be done.

\section*{Aetiology and Pathophysiology}
Continence depends on the pelvic floor muscles, the urethral sphincter, and normal bladder function. Stress incontinence occurs when the pelvic floor and sphincter are weakened, most often by childbirth, ageing, obesity, chronic coughing, or heavy lifting, so that rises in abdominal pressure overwhelm the sphincter. Urge incontinence results from an overactive detrusor muscle in the bladder wall, which contracts involuntarily at small volumes, often for no identifiable cause but sometimes after stroke, nerve disease, or bladder irritation. Factors that worsen all types include obesity, constipation, caffeine and alcohol, some medicines, urinary tract infections, and diabetes. In men, prostate surgery can damage the sphincter, while an enlarged prostate can cause urgency and overflow.

\section*{Clinical Features}
\begin{itemize}
    \item Leakage with coughing, sneezing, laughing, lifting, or exercise (stress incontinence)
    \item Sudden, strong urges to pass urine with leakage before reaching the toilet (urge incontinence)
    \item A combination of stress and urge symptoms (mixed incontinence)
    \item Passing urine very frequently, including at night
    \item A feeling of incomplete bladder emptying
    \item Constant slow leakage, which may indicate overflow incontinence
    \item The need to wear pads, and avoidance of social activities and exercise
    \item Effects on sleep, mood, intimacy, and quality of life
\end{itemize}

\section*{Differential Diagnosis}
Leakage may be due to conditions other than pelvic floor weakness or an overactive bladder.
\begin{itemize}
    \item \textbf{Urinary tract infection:} acute onset with burning and frequency
    \item \textbf{Overflow incontinence:} an overfull bladder leaking because of obstruction or weak bladder muscle
    \item \textbf{Transient causes:} medicines, excess fluid intake, constipation, and delirium
    \item \textbf{Neurological disease:} stroke, multiple sclerosis, Parkinson's disease, or spinal problems
    \item \textbf{Fistula:} an abnormal connection between the bladder and vagina, usually after surgery or childbirth in lower-resource settings
    \item \textbf{Prostate disease in men:} enlargement or cancer treatment
    \item \textbf{Functional problems:} difficulty reaching the toilet because of mobility or cognitive problems
\end{itemize}

\section*{When to Seek Medical Care}
Anyone troubled by urine leakage should discuss it with a doctor, because treatment is effective and early assessment gives the best results. Seek review particularly if leakage begins suddenly, is associated with pain or burning, or is accompanied by blood in the urine. In children, bedwetting and daytime leakage warrant assessment when they persist beyond the expected age or cause distress.

\textbf{Contact your local emergency services for:}
\begin{itemize}
    \item Sudden inability to pass urine with severe lower abdominal pain, which may indicate acute urinary retention
    \item Complete loss of bladder or bowel control with leg weakness or numbness, which may indicate spinal cord compression
    \item Passing large amounts of blood with difficulty passing urine
\end{itemize}

\section*{Diagnosis and Investigations}
\begin{itemize}
    \item \textbf{History:} type and timing of leakage, fluid intake, medicines, childbirth and surgery history
    \item \textbf{Bladder diary:} recording fluid intake and urine output over several days
    \item \textbf{Urinalysis:} to exclude infection, blood, or glucose
    \item \textbf{Examination:} assessment of the abdomen, pelvic floor muscle strength, and in women a vaginal examination
    \item \textbf{Post-void residual measurement:} ultrasound to check the bladder empties fully
    \item \textbf{Pad test and cough stress test:} simple ways to confirm leakage
    \item \textbf{Urodynamics:} specialised testing of bladder pressure and flow where the diagnosis is unclear or surgery is planned
    \item \textbf{Referral:} to a continence nurse, physiotherapist, urologist, or gynaecologist as needed
\end{itemize}

\section*{Management}
\begin{itemize}
    \item \textbf{Pelvic floor muscle training:} the first-line treatment for stress and mixed incontinence, requiring supervised teaching and regular practice
    \item \textbf{Bladder training:} scheduled voiding and techniques to resist urgency
    \item \textbf{Lifestyle measures:} weight loss, treating constipation, and reducing caffeine and alcohol
    \item \textbf{Medicines:} anticholinergic or beta-3 agonist drugs to calm an overactive bladder
    \item \textbf{Devices:} continence pessaries or urethral inserts in some women
    \item \textbf{Oestrogen therapy:} vaginal oestrogen for women with vaginal atrophy contributing to symptoms
    \item \textbf{Procedures:} botulinum toxin injections into the bladder, sacral nerve stimulation, or surgery such as a sling procedure for stress incontinence
    \item \textbf{Male-specific care:} treatment of prostate problems and post-prostatectomy rehabilitation
    \item \textbf{Intermittent catheterisation:} for overflow incontinence
    \item \textbf{Continence products and aids:} pads, protective garments, and aids that allow full participation in life
\end{itemize}

\section*{Complications}
\begin{itemize}
    \item Skin irritation, rashes, and infections from constant moisture
    \item Recurrent urinary tract infections
    \item Falls and fractures in older people hurrying to the toilet at night
    \item Sleep disturbance and fatigue
    \item Depression, anxiety, loss of confidence, and social isolation
    \item Effects on intimacy and relationships
    \item Reduced physical activity and its health consequences
    \item In some cases, worsening of kidney function from chronic obstruction
\end{itemize}

\section*{Prognosis and Outcomes}
Most people with incontinence improve significantly with treatment. Pelvic floor training cures or greatly improves stress incontinence in around two-thirds of women who practise it consistently, and bladder training and medicines control urgency in most. Weight loss of around five to ten per cent can markedly reduce leakage in overweight women. Surgery for stress incontinence has high success rates in appropriately selected women. Because continence problems are so treatable, the most important step is seeking assessment rather than accepting leakage as normal.

\section*{Prevention and Self-Care}
\begin{itemize}
    \item Practise pelvic floor exercises, especially during and after pregnancy
    \item Maintain a healthy weight
    \item Treat constipation and avoid straining
    \item Limit caffeine and alcohol, which irritate the bladder
    \item Drink enough fluid through the day without excessive amounts late at night
    \item See a continence physiotherapist for expert teaching
    \item Use the toilet when the urge is normal, and avoid habitual rushing
    \item Keep a bladder diary to guide treatment
    \item Talk to a doctor early rather than suffering in silence
    \item For carers: respond promptly and kindly to requests for toilet assistance
\end{itemize}

\section*{Sources and Further Reading}
The following reputable sources provide further information for healthcare students and patients seeking reliable, current advice about urinary incontinence.
\begin{itemize}
    \item Continence Foundation of Australia. \textit{Information about bladder and bowel health.} https://www.continence.org.au/
    \item Healthdirect Australia. \textit{Urinary incontinence.} https://www.healthdirect.gov.au/
    \item Jean Hailes for Women's Health. \textit{Bladder health and incontinence.} https://www.jeanhailes.org.au/
    \item NHS. \textit{Urinary incontinence.} https://www.nhs.uk/conditions/urinary-incontinence/
    \item Mayo Clinic. \textit{Urinary incontinence.} https://www.mayoclinic.org/
    \item International Continence Society. \textit{Patient information on continence.} https://www.ics.org/
\end{itemize}
